\chapter{S4：环境可靠性、复测与批量一致性}
\section{阶段定位与设备边界}
本阶段将 S1--S3 的关键指标转化为随时间、环境应力与批次变化的稳定性证据。单块 PCB 只适合研发诊断；批量一致性与工艺能力需要有代表性的样本和抽样计划。
\section{S4 测试项目总表}
\small
\begin{longtable}{C{18mm}p{46mm}p{61mm}p{22mm}}
\toprule 编号 & 测试项 & 测试目标／方法提示 & 级别／条件\\\midrule\endfirsthead
\toprule 编号 & 测试项 & 测试目标／方法提示 & 级别／条件\\\midrule\endhead
S4-00 & 样本量与可靠性计划审核 & 量产/研发样本区分、置信度和应力等级 & 前置/必做\\
S4-01 & 长期运行后零偏漂移 & 长时间留置、定期 S1 复测 & P0\\
S4-02 & 长期运行后 Allan 变化 & 对比 CT-12A 随老化/寿命的变化 & P1\\
S4-03 & 温度循环前后内参漂移 & 循环前后同一姿态/速率基准重复 & P0\\
S4-04 & 振动后内参及角度复测 & 应力前后残余误差差分 & P0/需设备\\
S4-05 & 机械冲击后参数和输出恢复 & 冲击前后变化、结构/焊点检查 & P1/需设备\\
S4-06 & 反复通断电和存储完整性 & NVM 参数、CRC、版本、启动重复性 & P0\\
S4-07 & 重复装夹 GR\&R & 操作员/夹具/台体贡献分解 & P1\\
S4-08 & 跨 PCB/批次一致性 & 零偏/矩阵/角度指标分布与离群点 & P0/多样本\\
S4-09 & 转台/参考设备交叉验证 & 不同参考条件下重复测量的一致性 & P1/需设备\\
S4-10 & 工艺能力与不良追踪 & 标定失败率、误判率、追溯覆盖率 & P0/量产\\
S4-11 & 产品放行与可靠性报告 & 评审准则、纠正措施、版本归档 & P0\\
\bottomrule
\end{longtable}
\normalsize
\section{逐项方案编写接口（后续版本）}
本阶段\textbf{测试编号已预留，详细实施方案尚未完成}。每个编号在正式执行前，应按 S1 章节的模板补齐：目的、仪器和参考不确定度、进入条件、工况矩阵、台架操作、原始数据字段、计算公式、合格判据、异常处理和结果归档。不得因本表列有测试名称就自动判定已经执行或通过。
\begin{warningbox}若只有一块 PCB，不能给出多批次一致性统计和 GR\&R 结论；此类条目只能先保留方法框架。\end{warningbox}
